% 问题重述与分析
\section{引言}
\subsection{问题背景}
生鲜商超是城市居民日常消费的重要渠道，其中蔬菜类商品因消费频率高、需求刚性强的特点，在商超经营中占据关键地位。然而，蔬菜类商品保鲜期短、品相变化快，商超需按日进行补货。并且，由于蔬菜品种与产地繁多，且进货交易时间集中在凌晨3:00至4:00，商超往往在信息不完全的情况下做出当日补货决策。同时，蔬菜类商品在需求侧呈现出销售量与时间关联关系，在供给侧中则表现出季节性，以及与销售空间限制的紧密联系。在此背景下，如何依据历史销售数据、批发价格和损耗率等信息，科学制定蔬菜类商品的补货、定价、销售组合策略，已成为生鲜商超经营管理的核心问题之一。

\subsection{问题重述}
本题基于某商超经销的6个蔬菜品类的商品信息（附件1）、2020年7月1日至2023年6月30日的销售流水明细数据（附件2）与批发价格数据（附件3），以及各商品近期的损耗率数据（附件4），解决以下四个问题：

\textbf{问题一：}要求基于附件1、2等数据来源，分析各蔬菜品类及单品在销售量上的分布规律，如时间分布特征（季节性、周期性等），并揭示相关关系，如不同品类之间、不同单品之间可能存在的互补或替代关系，为后续补货与定价决策提供依据。

\textbf{问题二：}要求基于附件1、2、3、4等数据来源，在品类层面上建立销售总量与成本加成定价之间的关系模型，并在此基础上确定未来一周（2023年7月1日至7日）日补货总量和定价策略，以最大化商超的总收益。

\textbf{问题三：}要求基于附件2等数据来源，在可售单品总数控制在27至33个，且各单品订购量满足最小陈列量2.5千克的基础下，根据2023年6月24日至30日的可售品种，制定7月1日的单品补货量和定价策略，在尽量满足市场需求前提下，使得商超收益最大。

\textbf{问题四：}要求从优化制定蔬菜类商品补货和定价决策的角度出发，识别当前数据体系的不足，提出补充采集数据的意见和理由，并论述这些数据如何帮助解决上述问题。

\subsection{问题分析}

\textbf{问题一：}属于统计规律刻画与关联结构识别问题。以附件1（6品类、251单品）与附件2（2020-07-01至2023-06-30销售流水，$T=1095$天）为输入，输出品类与单品的销量分布形态、时间特征及协同关联。三个统计事实决定建模路径：(1)品类日销量为连续正值、右偏长尾，用对数正态/伽马分布拟合，AIC/BIC选型；(2)单品日销量存在大量无销售日，采用“售出概率$+$正销量条件分布”的两阶段混合模型；(3)原始日销量相关混入共同日历与季节成分会产生伪相关\cite{yule1926}，须先剔除时间效应再构建关联网络；单品销售形态差异悬殊，以时序特征聚类分组。据此拆解为分布拟合、STL季节分解、去季节关联网络、时序特征聚类四个子模型。

\textbf{问题二：}属于品类层面的关系建模与随机优化决策问题。要求刻画销售总量与成本加成定价的关系（Log-Log弹性回归），并给出未来一周各品类的日补货总量与定价策略使收益最大。三个难点决定模型结构：(1)价格与销量的内生性：须先界定品类售价口径；(2)时序约束：最优加价率$m_c^*$只依赖品类弹性与损耗率、与批发价无关；(3)需求随机与损耗：按报童临界比$\kappa=(P^*-c')/P^*$在缺货与损耗间权衡。按“弹性回归$\to$最优定价$\to$去价格需求预测$\to$报童补货”四步建模，批发价以最近一周均价估计。

\textbf{问题三：}属于单品级联合组合优化问题。以2023-06-24至30可售品种为候选，同时决策单品是否上架、选哪档价、订多少货，满足可售总数27--33、单品最小陈列量2.5kg的耦合约束。与问题二相比：(1)决策粒度由品类下放到单品；(2)跨单品共享约束使问题无法按品类独立优化\cite{kok2009,vanryzin1999}。三个对策：(1)单品弹性借用品类弹性先验做Shrinkage估计\cite{james1961}；(2)售价离散为$K=11$档价格网格，转化为线性0-1 MIP保证全局最优\cite{wolsey1998}；(3)以缺货惩罚软约束量化“满足需求”，配合灵敏度给出权衡\cite{khouja1999}。模型复用问题一的聚类分组与问题二的弹性、损耗参数。

\textbf{问题四：}属于数据体系诊断与采集方案建议问题。从优化补货定价决策出发，识别现有数据不足，提出补充采集建议。基于建模过程定位五类缺失：(1)日销量聚合口径无法区分共同趋势与真实关联，购物篮关系不可识别；(2)售价口径混合正常销售与打折价，弹性向零衰减；(3)缺货日真实需求不可观测；(4)损耗率为静态均值、缺乏动态结构；(5)数据粒度为日级，缺少日内信息与外部竞品参照。据此从需求侧、市场侧、运营侧、供给侧与外部环境五个维度提出采集方案，每条按“补什么数据—改进哪个模型—改善哪个指标”的链条与前文挂钩。